\documentclass[11pt,a4paper]{article}

% ---------------------------------------------------------
% Encoding, Fonts, Language
% ---------------------------------------------------------
\usepackage[T1]{fontenc}
\usepackage[utf8]{inputenc}
\usepackage[english]{babel}
\usepackage{lmodern}

% ---------------------------------------------------------
% Math Packages
% ---------------------------------------------------------
\usepackage{amsmath, amssymb, amsfonts}
\usepackage{mathtools}
\usepackage{bm}            % bold math symbols
\usepackage{physics}       % \dv, \pdv, \ket, \bra, etc.
\usepackage{tensor}        % tensor indices
\usepackage{siunitx}       % units

% ---------------------------------------------------------
% Figures, Tables, Layout
% ---------------------------------------------------------
\usepackage{graphicx}
\usepackage{float}
\usepackage{booktabs}
\usepackage{caption}
\usepackage{subcaption}
\usepackage{geometry}
\geometry{margin=1in}

% ---------------------------------------------------------
% Bibliography
% ---------------------------------------------------------

\usepackage[numbers,sort&compress]{natbib}


% ---------------------------------------------------------
% Hyperref (must be loaded LAST except cleveref)
% ---------------------------------------------------------
\usepackage[
    colorlinks=true,
    linkcolor=blue,
    citecolor=blue,
    urlcolor=blue,
    pdfencoding=auto,
    psdextra
]{hyperref}

% ---------------------------------------------------------
% Cleveref (after hyperref)
% ---------------------------------------------------------
\usepackage{cleveref}

% ---------------------------------------------------------
% Fix for math in section titles (renamed to avoid conflicts)
% ---------------------------------------------------------
\newcommand{\mathsafe}[2]{\texorpdfstring{#1}{#2}}

% ---------------------------------------------------------
% Paragraph spacing
% ---------------------------------------------------------
\setlength{\parskip}{1em}
\setlength{\parindent}{0pt}

% ---------------------------------------------------------
% Title
% ---------------------------------------------------------
\title{\textbf{Boltzmann Hierarchy and Linear Perturbations in\\
Relaxation-Driven Cyclic Cosmology (RDCC)}}

\author{Michael Lehmann\\
Independent Researcher\\
\texttt{mi.lehmann@gmx.de}}

\date{May 2026}

\begin{document}

\maketitle


\begin{abstract}
Relaxation-Driven Cyclic Cosmology (RDCC) is a one-parameter, CPT-symmetric
framework in which all observable deviations from $\Lambda$CDM arise from the
infrared relaxation strength $\alpha_{\mathrm{IR}}$. Companions~I--VI established the
gravitational foundations, tensor phenomenology, relaxation dynamics, global CPT
structure, observational synthesis, and the microscopic evolution of $\alpha(a)$.
The resulting framework is dynamically complete and empirically predictive.

The present paper, Companion~VII, develops the full linear perturbation theory of
RDCC and constructs the Boltzmann hierarchy for scalar, vector, and tensor modes
in both the visible and mirror sectors. We derive the modified Einstein equations,
the coupled continuity and Euler equations, the photon and neutrino Boltzmann
hierarchies, and the tensor-mode evolution including the RDCC modulation
mechanism. The relaxation parameter enters the perturbation system through a
modified Poisson equation, an effective gravitational coupling, and a small
inter-sector energy-exchange term.

We show that the freezing of $\alpha(a)$ established in Companion~VI implies that
all perturbative corrections are controlled by the single constant parameter
$\alpha_{\mathrm{IR}}$. This leads to a unified one-parameter modification of the CMB
anisotropies, matter power spectrum, and stochastic gravitational-wave background.
We provide the complete set of equations required for numerical implementation in
Boltzmann codes such as CLASS or CAMB, thereby enabling a full confrontation of
RDCC with precision cosmological data.

Companion~VII completes the predictive architecture of RDCC and establishes the
framework as a fully testable alternative to inflationary $\Lambda$CDM.
\end{abstract}

\newpage

\tableofcontents

\newpage

\section{Introduction}

Relaxation-Driven Cyclic Cosmology (RDCC) is a minimal, CPT-symmetric framework in
which all observable deviations from $\Lambda$CDM arise from a single infrared parameter:
the relaxation strength $\alpha_{\mathrm{IR}} \sim 10^{-2}$. Companions~I–VI developed the geometric
foundations, tensor phenomenology, relaxation dynamics, global CPT structure, observational
synthesis, and the microscopic evolution of $\alpha(a)$. The resulting picture is simple and
structurally robust: the relaxation parameter grows rapidly near the CPT-symmetric bounce,
freezes by $a \sim 10^{-4}$, and remains constant throughout all observable epochs.

This early freeze-out implies that the RDCC prediction vector


\[
    (n_s,\ \Delta N_{\mathrm{eff}},\ f\sigma_8,\ \Omega_{\mathrm{GW}})
\]


is controlled entirely by the single infrared value $\alpha_{\mathrm{IR}}$. Companion~V demonstrated
that this one-parameter structure is compatible with current CMB, LSS, and GW data, and
Companion~VI provided the dynamical justification for the freeze-out of $\alpha(a)$.

The purpose of the present paper, Companion~VII, is to develop the full linear perturbation
theory of RDCC and construct the Boltzmann hierarchy for scalar, vector, and tensor modes.
This is the final theoretical step required to make RDCC a fully testable cosmological
framework. The Boltzmann hierarchy determines the evolution of:
\begin{itemize}
    \item photon temperature and polarization multipoles,
    \item neutrino multipoles,
    \item baryon and CDM density and velocity perturbations,
    \item mirror-sector analogues of all above quantities,
    \item and the tensor-mode spectrum including the RDCC modulation mechanism.
\end{itemize}

The relaxation parameter enters the perturbation system in three ways:
\begin{enumerate}
    \item a modified Poisson equation with an effective gravitational coupling,
    \item a small inter-sector energy-exchange term proportional to $\alpha_{\mathrm{IR}}$,
    \item and a tensor-sector modulation frequency proportional to $\alpha_{\mathrm{IR}}$.
\end{enumerate}

Because $\alpha(a)$ is frozen during all observable epochs, these modifications are
time-independent and controlled entirely by $\alpha_{\mathrm{IR}}$. This makes RDCC one of the
most predictive alternatives to inflationary $\Lambda$CDM: all perturbative signatures arise
from a single microscopic parameter.

The remainder of this companion paper is organized as follows. Section~2 develops the
background cosmology and the effective gravitational coupling. Section~3 derives the scalar
perturbation equations. Section~4 constructs the photon and neutrino Boltzmann hierarchies.
Section~5 develops the mirror-sector perturbations. Section~6 presents the tensor-mode
evolution including the RDCC modulation mechanism. Section~7 assembles the full set of
equations required for numerical implementation in Boltzmann codes. Section~8 provides the
conclusion and outlines the path toward a full RDCC likelihood analysis.


\section{Background Cosmology and Effective Gravitational Coupling}

The Boltzmann hierarchy evolves perturbations on top of a homogeneous and isotropic
background. In RDCC, the background cosmology differs from $\Lambda$CDM in two
structural ways:

\begin{enumerate}
    \item the total energy density receives contributions from both the visible and mirror
    sectors, which share a common CPT-symmetric origin but evolve with different
    temperatures,
    \item the effective gravitational coupling is modified by the infrared relaxation strength
    $\alpha_{\mathrm{IR}}$, reflecting the residual inter-sector decoherence after the bounce.
\end{enumerate}

Because $\alpha(a)$ freezes extremely early (Companion~VI), all background modifications
are constant during observable epochs and controlled entirely by the single infrared parameter
$\alpha_{\mathrm{IR}}$. This section develops the background equations required for the perturbation
analysis.

\subsection{Two-Sector Background}

The total energy density is the sum of the visible and mirror contributions:
\begin{equation}
    \rho_{\mathrm{tot}}(a) = \rho_{+}(a) + \rho_{-}(a),
\end{equation}
with analogous expressions for pressure.

The mirror sector contains mirror photons, mirror neutrinos, mirror baryons, and (optionally)
mirror CDM. Their temperatures are related by the relaxation-induced ratio
\begin{equation}
    \frac{T_{-}}{T_{+}} = x, \qquad x \simeq 0.5 + \mathcal{O}(\alpha_{\mathrm{IR}}),
\end{equation}
as derived in Companion~III.

The Friedmann equation becomes
\begin{equation}
    H^{2}(a)
    = \frac{8\pi G}{3}\,\bigl[\rho_{+}(a) + \rho_{-}(a)\bigr].
\end{equation}

\subsection{Effective Gravitational Coupling}

The dimension-six inter-sector operator modifies the gravitational response of visible-sector
perturbations. As shown in Companions~III and~IV, the effective Newton constant is
\begin{equation}
    G_{\mathrm{eff}} = G\bigl(1 + c_{g}\,\alpha_{\mathrm{IR}}\bigr),
\end{equation}
where $c_{g}$ is an order-one coefficient determined by the inter-sector coupling.

This modification enters the Poisson equation and therefore affects:
\begin{itemize}
    \item the growth of structure,
    \item the CMB lensing potential,
    \item the matter power spectrum,
    \item and the scalar Boltzmann hierarchy.
\end{itemize}

Because $\alpha(a)$ is frozen, $G_{\mathrm{eff}}$ is constant during all observable epochs.

\subsection{Modified Expansion History}

The relaxation dynamics induce a small correction to the Hubble parameter:
\begin{equation}
    H^{2}(a) = H^{2}_{\Lambda\mathrm{CDM}}(a)\,\bigl[1 + \delta_{H}(a)\bigr],
\end{equation}
with
\begin{equation}
    \delta_{H}(a) \sim \alpha_{\mathrm{IR}}\,f(a),
\end{equation}
where $f(a)$ is a slowly varying function determined by the relative abundance of mirror
radiation and mirror matter.

To leading order in $\alpha_{\mathrm{IR}}$,
\begin{equation}
    H(a) = H_{\Lambda\mathrm{CDM}}(a)\,\bigl[1 + \tfrac{1}{2}\delta_{H}(a)\bigr].
\end{equation}

This correction is small ($\sim 1\%$) but must be included for consistency in the Boltzmann
hierarchy.

\subsection{Background Quantities Required for Perturbations}

The Boltzmann hierarchy requires the following background functions:
\begin{itemize}
    \item $H(a)$ and $H'(a)$,
    \item $\rho_{+}(a)$, $p_{+}(a)$,
    \item $\rho_{-}(a)$, $p_{-}(a)$,
    \item the temperature ratio $x = T_{-}/T_{+}$,
    \item the effective Newton constant $G_{\mathrm{eff}}$,
    \item the fractional energy densities $\Omega_{i}(a)$ for all species.
\end{itemize}

All of these quantities are determined by:
\begin{enumerate}
    \item the standard $\Lambda$CDM parameters,
    \item the temperature ratio $x$,
    \item and the single RDCC parameter $\alpha_{\mathrm{IR}}$.
\end{enumerate}

\subsection{Summary}

The background cosmology of RDCC differs from $\Lambda$CDM only through:
\begin{itemize}
    \item the presence of a mirror sector with temperature ratio $x$,
    \item a modified gravitational coupling $G_{\mathrm{eff}}$,
    \item and a small correction to the expansion history.
\end{itemize}

These modifications are constant during all observable epochs and controlled entirely by
$\alpha_{\mathrm{IR}}$. The next section develops the scalar perturbation equations on this
background.

\section{Scalar Perturbations}

Scalar perturbations form the backbone of the Boltzmann hierarchy: they determine the CMB
temperature and polarization anisotropies, the matter power spectrum, and the growth of
structure. In RDCC, the scalar sector differs from $\Lambda$CDM in three structural ways:

\begin{enumerate}
    \item the Poisson equation contains the effective gravitational coupling $G_{\mathrm{eff}}$,
    \item the visible and mirror sectors contribute separately to the gravitational potentials,
    \item a small inter-sector relaxation term proportional to $\alpha_{\mathrm{IR}}$ appears in the
    continuity and Euler equations.
\end{enumerate}

Because $\alpha(a)$ freezes extremely early (Companion~VI), all modifications are constant
during observable epochs and controlled entirely by the single infrared parameter
$\alpha_{\mathrm{IR}}$. We work in conformal Newtonian gauge throughout:
\begin{equation}
    ds^{2} = a^{2}(\tau)\,\bigl[-(1+2\Phi)\,d\tau^{2} + (1-2\Psi)\,dx^{2}\bigr].
\end{equation}

\subsection{Modified Einstein Equations}

The scalar Einstein equations receive contributions from both sectors. The perturbed $(0,0)$
equation becomes
\begin{equation}
    k^{2}\Psi + 3\mathcal{H}(\Psi' + \mathcal{H}\Phi)
    = 4\pi G_{\mathrm{eff}}\,a^{2}\,(\delta\rho_{+} + \delta\rho_{-}),
\end{equation}
where $\mathcal{H} = a'/a$ is the conformal Hubble parameter.

The effective gravitational coupling is
\begin{equation}
    G_{\mathrm{eff}} = G\bigl(1 + c_{g}\,\alpha_{\mathrm{IR}}\bigr),
\end{equation}
as derived in Companions~III–IV and justified dynamically in Companion~VI.

The $(0,i)$ equation becomes
\begin{equation}
    k^{2}(\Psi' + \mathcal{H}\Phi)
    = 4\pi G_{\mathrm{eff}}\,a^{2}\bigl[(\rho_{+}+p_{+})\theta_{+}
    + (\rho_{-}+p_{-})\theta_{-}\bigr].
\end{equation}

The anisotropic-stress equation is unchanged:
\begin{equation}
    k^{2}(\Phi - \Psi)
    = 12\pi G\,a^{2}\bigl[(\rho_{+}+p_{+})\sigma_{+}
    + (\rho_{-}+p_{-})\sigma_{-}\bigr].
\end{equation}

\subsection{Continuity Equations}

The visible-sector continuity equation receives a small correction from the relaxation dynamics:
\begin{equation}
    \delta'_{+}
    = - (1+w_{+})(\theta_{+} - 3\Psi')
      - 3\mathcal{H}(c_{s,+}^{2} - w_{+})\delta_{+}
      + \alpha_{\mathrm{IR}}\mathcal{H}(\delta_{-} - \delta_{+}).
\end{equation}

The mirror-sector continuity equation is
\begin{equation}
    \delta'_{-}
    = - (1+w_{-})(\theta_{-} - 3\Psi')
      - 3\mathcal{H}(c_{s,-}^{2} - w_{-})\delta_{-}
      - \alpha_{\mathrm{IR}}\mathcal{H}(\delta_{-} - \delta_{+}).
\end{equation}

The last terms in both equations represent the inter-sector relaxation. Because
$\alpha_{\mathrm{IR}} \sim 10^{-2}$, these terms are small but non-negligible.

\subsection{Euler Equations}

The Euler equations for the two sectors are
\begin{equation}
    \theta'_{+}
    = - \mathcal{H}(1 - 3c_{s,+}^{2})\theta_{+}
      + \frac{c_{s,+}^{2}}{1+w_{+}}\,k^{2}\delta_{+}
      + k^{2}\Phi
      + \alpha_{\mathrm{IR}}\mathcal{H}(\theta_{-} - \theta_{+}),
\end{equation}
\begin{equation}
    \theta'_{-}
    = - \mathcal{H}(1 - 3c_{s,-}^{2})\theta_{-}
      + \frac{c_{s,-}^{2}}{1+w_{-}}\,k^{2}\delta_{-}
      + k^{2}\Phi
      - \alpha_{\mathrm{IR}}\mathcal{H}(\theta_{-} - \theta_{+}).
\end{equation}

The relaxation term damps the velocity difference between the sectors, reflecting the residual
inter-sector decoherence encoded in $\alpha_{\mathrm{IR}}$.

\subsection{Adiabatic Initial Conditions}

Because the two sectors share a common origin at the CPT-symmetric bounce, adiabatic initial
conditions require
\begin{equation}
    \frac{\delta_{-}}{1+w_{-}}
    = \frac{\delta_{+}}{1+w_{+}},
    \qquad
    \theta_{-} = \theta_{+}.
\end{equation}

The relaxation term preserves adiabaticity at early times.

\subsection{Summary}

Scalar perturbations in RDCC differ from $\Lambda$CDM through:
\begin{itemize}
    \item a modified Poisson equation with $G_{\mathrm{eff}}$,
    \item separate visible and mirror contributions to the gravitational potentials,
    \item a small relaxation term proportional to $\alpha_{\mathrm{IR}}$ in the continuity and Euler
    equations.
\end{itemize}

These modifications are constant in time and controlled entirely by the single parameter
$\alpha_{\mathrm{IR}}$. The next section develops the photon and neutrino Boltzmann hierarchies
on this background.

\section{Photon and Neutrino Boltzmann Hierarchies}

The photon and neutrino Boltzmann hierarchies determine the evolution of the CMB
temperature and polarization anisotropies and the free-streaming behavior of relativistic
species. In RDCC, the structure of the Boltzmann hierarchy is identical to $\Lambda$CDM
except for three modifications:

\begin{enumerate}
    \item the gravitational potentials $(\Phi,\Psi)$ satisfy the modified Einstein equations of
    Sec.~3 with the effective coupling $G_{\mathrm{eff}}$,
    \item the mirror sector contributes to the gravitational potentials,
    \item mirror photons and mirror neutrinos obey parallel hierarchies with temperature
    ratio $x = T_{-}/T_{+}$.
\end{enumerate}

The relaxation parameter $\alpha_{\mathrm{IR}}$ does \emph{not} enter the photon or neutrino hierarchies
directly; it affects them only through the modified gravitational potentials. Because
$\alpha(a)$ is frozen during all observable epochs, these modifications are constant in time.

We work in conformal Newtonian gauge and follow the standard multipole expansion of the
photon and neutrino distribution functions.

\subsection{Photon Temperature Hierarchy}

The photon temperature perturbation is expanded as
\begin{equation}
    \Theta(\mathbf{k},\hat{n},\tau)
    = \sum_{\ell=0}^{\infty} (-i)^{\ell} (2\ell+1)\,\Theta_{\ell}(\mathbf{k},\tau)\,P_{\ell}(\mu),
\end{equation}
with $\mu = \hat{k}\!\cdot\!\hat{n}$.

The Boltzmann hierarchy is
\begin{align}
    \Theta_{0}' &= -k\Theta_{1} - \Phi', \\
    \Theta_{1}' &= \frac{k}{3}\Theta_{0} - \frac{2k}{3}\Theta_{2} + k\Psi - \dot{\tau}\,(\Theta_{1} - v_{b}), \\
    \Theta_{2}' &= \frac{2k}{5}\Theta_{1} - \frac{3k}{5}\Theta_{3}
                  + \dot{\tau}\left(\frac{1}{10}\Pi - \Theta_{2}\right), \\
    \Theta_{\ell}' &= \frac{k}{2\ell+1}\bigl[\ell\,\Theta_{\ell-1} - (\ell+1)\Theta_{\ell+1}\bigr]
                     - \dot{\tau}\,\Theta_{\ell}, \qquad \ell \ge 3.
\end{align}

The polarization source is
\begin{equation}
    \Pi = \Theta_{2} + \Theta_{2}^{(P)},
\end{equation}
with $\Theta_{2}^{(P)}$ the polarization quadrupole.

The only RDCC modification is that the gravitational potentials $(\Phi,\Psi)$ satisfy the
modified Einstein equations of Sec.~3.

\subsection{Photon Polarization Hierarchy}

The polarization multipoles obey
\begin{align}
    \Theta_{0}^{(P)\prime} &= -k\Theta_{1}^{(P)} - \dot{\tau}\,\Theta_{0}^{(P)}, \\
    \Theta_{1}^{(P)\prime} &= \frac{k}{3}\Theta_{0}^{(P)} - \frac{2k}{3}\Theta_{2}^{(P)}
                             - \dot{\tau}\,\Theta_{1}^{(P)}, \\
    \Theta_{2}^{(P)\prime} &= \frac{2k}{5}\Theta_{1}^{(P)} - \frac{3k}{5}\Theta_{3}^{(P)}
                             + \dot{\tau}\left(\frac{1}{10}\Pi - \Theta_{2}^{(P)}\right), \\
    \Theta_{\ell}^{(P)\prime} &=
    \frac{k}{2\ell+1}\bigl[\ell\,\Theta_{\ell-1}^{(P)} - (\ell+1)\Theta_{\ell+1}^{(P)}\bigr]
    - \dot{\tau}\,\Theta_{\ell}^{(P)}, \qquad \ell \ge 3.
\end{align}

Again, the only RDCC modification enters through the gravitational potentials.

\subsection{Neutrino Hierarchy}

Massless neutrinos free-stream and therefore require a full Boltzmann hierarchy:
\begin{align}
    F_{0}' &= -kF_{1} - \Phi', \\
    F_{1}' &= \frac{k}{3}F_{0} - \frac{2k}{3}F_{2} + k\Psi, \\
    F_{2}' &= \frac{2k}{5}F_{1} - \frac{3k}{5}F_{3}
             + \frac{4}{15}(\Phi' + \Psi'), \\
    F_{\ell}' &= \frac{k}{2\ell+1}\bigl[\ell\,F_{\ell-1} - (\ell+1)F_{\ell+1}\bigr],
               \qquad \ell \ge 3.
\end{align}

The RDCC modification again enters only through the gravitational potentials.

\subsection{Mirror-Photon and Mirror-Neutrino Hierarchies}

The mirror sector contains mirror photons and mirror neutrinos with temperature
\begin{equation}
    T_{-} = x\,T_{+}, \qquad x \simeq 0.5 + \mathcal{O}(\alpha_{\mathrm{IR}}).
\end{equation}

Their Boltzmann hierarchies are identical to the visible-sector ones, with the replacements


\[
    T_{+} \rightarrow T_{-}, \qquad
    \dot{\tau} \rightarrow \dot{\tau}_{-}, \qquad
    \Theta_{\ell} \rightarrow \Theta_{\ell,-},
\]


and with gravitational potentials sourced by both sectors.

\subsection{Summary}

The photon and neutrino Boltzmann hierarchies in RDCC are structurally identical to
$\Lambda$CDM, with the following modifications:

\begin{itemize}
    \item the gravitational potentials $(\Phi,\Psi)$ obey the modified Einstein equations,
    \item the mirror sector contributes to the potentials,
    \item the effective gravitational coupling $G_{\mathrm{eff}}$ modifies the Poisson equation,
    \item mirror photons and mirror neutrinos obey parallel hierarchies with temperature
    ratio $x$.
\end{itemize}

The next section develops the perturbation equations for mirror baryons and mirror CDM,
completing the scalar sector of RDCC.

\section{Mirror-Sector Perturbations}

The mirror sector contains mirror photons, mirror neutrinos, mirror baryons, and (optionally)
mirror CDM. Their perturbations obey the same dynamical equations as their visible-sector
counterparts, with two structural differences:

\begin{enumerate}
    \item the mirror temperature is lower by the relaxation-induced ratio
    

\[
        x = \frac{T_{-}}{T_{+}}, \qquad x \simeq 0.5 + \mathcal{O}(\alpha_{\mathrm{IR}}),
    \]


    \item the gravitational potentials $(\Phi,\Psi)$ are sourced jointly by the visible and mirror
    sectors.
\end{enumerate}

The relaxation parameter $\alpha_{\mathrm{IR}}$ enters the mirror-sector perturbations only through
the small inter-sector coupling terms derived in Sec.~3. Because $\alpha(a)$ is frozen during all
observable epochs, these terms are constant in time and controlled entirely by $\alpha_{\mathrm{IR}}$.

\subsection{Mirror Baryons}

Mirror baryons obey the continuity and Euler equations
\begin{align}
    \delta'_{b,-}
    &= -(\theta_{b,-} - 3\Psi')
       - 3\mathcal{H}(c_{s,b,-}^{2} - w_{b,-})\,\delta_{b,-}
       + \alpha_{\mathrm{IR}}\mathcal{H}(\delta_{b,+} - \delta_{b,-}), \\
    \theta'_{b,-}
    &= -\mathcal{H}(1 - 3c_{s,b,-}^{2})\,\theta_{b,-}
       + \frac{c_{s,b,-}^{2}}{1+w_{b,-}}\,k^{2}\delta_{b,-}
       + k^{2}\Phi
       - \alpha_{\mathrm{IR}}\mathcal{H}(\theta_{b,-} - \theta_{b,+}).
\end{align}

The Thomson scattering rate in the mirror sector is
\begin{equation}
    \dot{\tau}_{-} = a\,n_{e,-}\,\sigma_{T},
\end{equation}
with $n_{e,-}$ determined by the mirror recombination history. Because $T_{-} = xT_{+}$ with
$x < 1$, mirror recombination occurs earlier than visible recombination.

\subsection{Mirror CDM}

Mirror CDM (if present) obeys
\begin{align}
    \delta'_{c,-}
    &= -\theta_{c,-} + 3\Psi'
       + \alpha_{\mathrm{IR}}\mathcal{H}(\delta_{c,+} - \delta_{c,-}), \\
    \theta'_{c,-}
    &= -\mathcal{H}\,\theta_{c,-}
       + k^{2}\Phi
       - \alpha_{\mathrm{IR}}\mathcal{H}(\theta_{c,-} - \theta_{c,+}).
\end{align}

The relaxation term damps the velocity difference between visible and mirror CDM, reflecting
the residual inter-sector decoherence encoded in $\alpha_{\mathrm{IR}}$.

\subsection{Mirror Photons}

Mirror photons obey the same Boltzmann hierarchy as visible photons:
\begin{align}
    \Theta'_{0,-} &= -k\Theta_{1,-} - \Phi', \\
    \Theta'_{1,-} &= \frac{k}{3}\Theta_{0,-} - \frac{2k}{3}\Theta_{2,-}
                    + k\Psi - \dot{\tau}_{-}(\Theta_{1,-} - v_{b,-}), \\
    \Theta'_{2,-} &= \frac{2k}{5}\Theta_{1,-} - \frac{3k}{5}\Theta_{3,-}
                    + \dot{\tau}_{-}\left(\frac{1}{10}\Pi_{-} - \Theta_{2,-}\right), \\
    \Theta'_{\ell,-} &= \frac{k}{2\ell+1}\bigl[\ell\,\Theta_{\ell-1,-}
                      - (\ell+1)\Theta_{\ell+1,-}\bigr]
                      - \dot{\tau}_{-}\Theta_{\ell,-}, \qquad \ell \ge 3.
\end{align}

The only differences relative to the visible sector are:

\begin{itemize}
    \item the temperature ratio $x$,
    \item the mirror Thomson rate $\dot{\tau}_{-}$,
    \item and the gravitational potentials sourced by both sectors.
\end{itemize}

\subsection{Mirror Neutrinos}

Mirror neutrinos free-stream and obey the hierarchy
\begin{align}
    F'_{0,-} &= -kF_{1,-} - \Phi', \\
    F'_{1,-} &= \frac{k}{3}F_{0,-} - \frac{2k}{3}F_{2,-} + k\Psi, \\
    F'_{2,-} &= \frac{2k}{5}F_{1,-} - \frac{3k}{5}F_{3,-}
               + \frac{4}{15}(\Phi' + \Psi'), \\
    F'_{\ell,-} &= \frac{k}{2\ell+1}\bigl[\ell\,F_{\ell-1,-}
                 - (\ell+1)F_{\ell+1,-}\bigr], \qquad \ell \ge 3.
\end{align}

The mirror neutrino temperature is
\begin{equation}
    T_{\nu,-} = x\,T_{\nu,+}.
\end{equation}

\subsection{Gravitational Coupling Between Sectors}

The gravitational potentials are sourced by both sectors:
\begin{equation}
    k^{2}\Phi
    = 4\pi G_{\mathrm{eff}}\,a^{2}\,(\delta\rho_{+} + \delta\rho_{-})
      + 3\mathcal{H}(\Phi' + \mathcal{H}\Psi).
\end{equation}

This coupling is the dominant interaction between the sectors. The relaxation terms
proportional to $\alpha_{\mathrm{IR}}$ are subdominant but necessary for consistency.

\subsection{Summary}

Mirror-sector perturbations in RDCC:
\begin{itemize}
    \item obey the same dynamical equations as visible-sector perturbations,
    \item have different temperatures and recombination histories,
    \item contribute to the gravitational potentials,
    \item and are weakly coupled to the visible sector through $\alpha_{\mathrm{IR}}$.
\end{itemize}

The next section develops the tensor-mode evolution, including the RDCC modulation
mechanism derived in Companion~II.

\section{Tensor Perturbations and RDCC Modulation}

Tensor perturbations describe gravitational waves propagating on the FLRW background.
In RDCC, the tensor sector contains a distinctive and predictive signature: a log-periodic
modulation of the tensor amplitude with frequency proportional to the infrared relaxation
strength $\alpha_{\mathrm{IR}}$. This modulation was derived in Companion~II, interpreted in
Companions~III–IV, and dynamically justified in Companion~VI. The purpose of this section
is to embed that structure into the full Boltzmann hierarchy.

\subsection{Tensor Evolution Equation}

In conformal Newtonian gauge, the tensor perturbation $h_{ij}$ satisfies
\begin{equation}
    h'' + 2\mathcal{H}h' + k^{2}h
    = 16\pi G_{\mathrm{eff}}\,a^{2}\,\pi^{(T)},
\end{equation}
where $\pi^{(T)}$ is the tensor anisotropic stress.

The RDCC modification enters through the effective Newton constant
\begin{equation}
    G_{\mathrm{eff}} = G\bigl(1 + c_{g}\,\alpha_{\mathrm{IR}}\bigr).
\end{equation}

Because $\alpha(a)$ is frozen during all observable epochs, the modification is constant in
time.

\subsection{Tensor Anisotropic Stress}

The tensor anisotropic stress receives contributions from:
\begin{itemize}
    \item visible neutrinos,
    \item mirror neutrinos,
    \item and any additional free-streaming relativistic species.
\end{itemize}

The total stress is
\begin{equation}
    \pi^{(T)} = \pi^{(T)}_{\nu,+} + \pi^{(T)}_{\nu,-}.
\end{equation}

Each component obeys the standard hierarchy
\begin{align}
    F^{(T)\prime}_{2} &= -\frac{8}{15}h' - \frac{3k}{5}F^{(T)}_{3}, \\
    F^{(T)\prime}_{\ell} &=
    \frac{k}{2\ell+1}\bigl[\ell\,F^{(T)}_{\ell-1}
    - (\ell+1)F^{(T)}_{\ell+1}\bigr], \qquad \ell \ge 3,
\end{align}
with temperature ratio $T_{\nu,-} = xT_{\nu,+}$ for the mirror sector.

\subsection{RDCC Tensor Modulation}

Companion~II showed that the dimension-six operator induces a modulation of the tensor
amplitude:
\begin{equation}
    h(k,\tau)
    = h_{0}(k,\tau)\left[1 + \alpha_{\mathrm{IR}}
      \cos\!\left(\omega_{\mathrm{mod}}\ln\frac{k}{k_{\ast}}\right)\right],
\end{equation}
where
\begin{equation}
    \omega_{\mathrm{mod}} = c_{w}\,\alpha_{\mathrm{IR}},
\end{equation}
and $c_{w}$ is an order-one coefficient determined in Companion~II.

This modulation is a direct consequence of the relaxation dynamics and the bipartite
Hilbert-space structure of RDCC.

\subsection{Embedding the Modulation into the Boltzmann Hierarchy}

To incorporate the modulation into the Boltzmann evolution, we write the tensor mode as
\begin{equation}
    h(k,\tau) = h_{\mathrm{RDCC}}(k)\,h_{\mathrm{std}}(k,\tau),
\end{equation}
where
\begin{equation}
    h_{\mathrm{RDCC}}(k)
    = 1 + \alpha_{\mathrm{IR}}
      \cos\!\left(\omega_{\mathrm{mod}}\ln\frac{k}{k_{\ast}}\right)
\end{equation}
is the primordial modulation factor, and $h_{\mathrm{std}}(k,\tau)$ obeys the standard tensor
evolution equation with $G_{\mathrm{eff}}$.

This separation ensures that:
\begin{itemize}
    \item the modulation affects only the initial conditions,
    \item the time evolution remains standard (up to $G_{\mathrm{eff}}$),
    \item and the RDCC signature propagates cleanly into the observable spectrum.
\end{itemize}

\subsection{Modified Initial Conditions}

The initial tensor power spectrum becomes
\begin{equation}
    P_{h}(k)
    = P_{h}^{(0)}(k)\left[1 + \alpha_{\mathrm{IR}}
      \cos\!\left(\omega_{\mathrm{mod}}\ln\frac{k}{k_{\ast}}\right)\right]^{2}.
\end{equation}

To linear order in $\alpha_{\mathrm{IR}}$,
\begin{equation}
    P_{h}(k)
    \simeq P_{h}^{(0)}(k)\left[1 + 2\alpha_{\mathrm{IR}}
      \cos\!\left(\omega_{\mathrm{mod}}\ln\frac{k}{k_{\ast}}\right)\right].
\end{equation}

This is the form implemented in Boltzmann codes.

\subsection{Observable Consequences}

The RDCC tensor modulation produces:
\begin{itemize}
    \item a log-periodic oscillation in the stochastic gravitational-wave background,
    \item a characteristic modulation in the tensor transfer function,
    \item a distinctive signature in the CMB $B$-mode spectrum,
    \item and a frequency-dependent oscillation in $\Omega_{\mathrm{GW}}(f)$.
\end{itemize}

These signatures are controlled entirely by $\alpha_{\mathrm{IR}}$ and therefore provide clean,
parameter-free tests of RDCC.

\subsection{Summary}

Tensor perturbations in RDCC differ from $\Lambda$CDM through:
\begin{itemize}
    \item a modified gravitational coupling $G_{\mathrm{eff}}$,
    \item contributions from mirror neutrinos to the anisotropic stress,
    \item and a primordial log-periodic modulation with frequency
    $\omega_{\mathrm{mod}} \propto \alpha_{\mathrm{IR}}$.
\end{itemize}

The next section assembles the full set of equations required for numerical implementation
in Boltzmann codes such as CLASS or CAMB.

\section{Full Equation Set for Boltzmann Implementation}

This section assembles the complete set of equations required to implement RDCC in a
Boltzmann code such as CLASS or CAMB. The structure follows the standard
$\Lambda$CDM architecture, with RDCC modifications highlighted explicitly. All equations
are written in conformal Newtonian gauge.

Because the relaxation parameter freezes extremely early (Companion~VI), all RDCC
modifications are constant in time and controlled entirely by the single infrared parameter
$\alpha_{\mathrm{IR}}$.

\subsection{Background Evolution}

\paragraph{Friedmann equation.}
\begin{equation}
    H^{2}(a)
    = \frac{8\pi G}{3}\,\bigl[\rho_{+}(a) + \rho_{-}(a)\bigr]\,
      \bigl[1 + \delta_{H}(a)\bigr],
\end{equation}
with
\begin{equation}
    \delta_{H}(a) \sim \alpha_{\mathrm{IR}}\,f(a).
\end{equation}

\paragraph{Effective Newton constant.}
\begin{equation}
    G_{\mathrm{eff}} = G\bigl(1 + c_{g}\,\alpha_{\mathrm{IR}}\bigr).
\end{equation}

\paragraph{Temperature ratio.}
\begin{equation}
    T_{-} = x\,T_{+}, \qquad x \simeq 0.5 + \mathcal{O}(\alpha_{\mathrm{IR}}).
\end{equation}

\paragraph{Energy densities.}
\begin{equation}
    \rho_{\mathrm{tot}} = \rho_{+} + \rho_{-}, \qquad
    p_{\mathrm{tot}} = p_{+} + p_{-}.
\end{equation}

\subsection{Scalar Einstein Equations}

\paragraph{Poisson equation.}
\begin{equation}
    k^{2}\Psi + 3\mathcal{H}(\Psi' + \mathcal{H}\Phi)
    = 4\pi G_{\mathrm{eff}}\,a^{2}\,(\delta\rho_{+} + \delta\rho_{-}).
\end{equation}

\paragraph{Momentum constraint.}
\begin{equation}
    k^{2}(\Psi' + \mathcal{H}\Phi)
    = 4\pi G_{\mathrm{eff}}\,a^{2}\bigl[(\rho_{+}+p_{+})\theta_{+}
    + (\rho_{-}+p_{-})\theta_{-}\bigr].
\end{equation}

\paragraph{Anisotropic stress.}
\begin{equation}
    k^{2}(\Phi - \Psi)
    = 12\pi G\,a^{2}\bigl[(\rho_{+}+p_{+})\sigma_{+}
    + (\rho_{-}+p_{-})\sigma_{-}\bigr].
\end{equation}

\subsection{Scalar Continuity and Euler Equations}

\paragraph{Visible sector.}
\begin{align}
    \delta'_{+}
    &= - (1+w_{+})(\theta_{+} - 3\Psi')
       - 3\mathcal{H}(c_{s,+}^{2} - w_{+})\delta_{+}
       + \alpha_{\mathrm{IR}}\mathcal{H}(\delta_{-} - \delta_{+}), \\
    \theta'_{+}
    &= - \mathcal{H}(1 - 3c_{s,+}^{2})\theta_{+}
       + \frac{c_{s,+}^{2}}{1+w_{+}}\,k^{2}\delta_{+}
       + k^{2}\Phi
       + \alpha_{\mathrm{IR}}\mathcal{H}(\theta_{-} - \theta_{+}).
\end{align}

\paragraph{Mirror sector.}
\begin{align}
    \delta'_{-}
    &= - (1+w_{-})(\theta_{-} - 3\Psi')
       - 3\mathcal{H}(c_{s,-}^{2} - w_{-})\delta_{-}
       - \alpha_{\mathrm{IR}}\mathcal{H}(\delta_{-} - \delta_{+}), \\
    \theta'_{-}
    &= - \mathcal{H}(1 - 3c_{s,-}^{2})\theta_{-}
       + \frac{c_{s,-}^{2}}{1+w_{-}}\,k^{2}\delta_{-}
       + k^{2}\Phi
       - \alpha_{\mathrm{IR}}\mathcal{H}(\theta_{-} - \theta_{+}).
\end{align}

\subsection{Photon and Neutrino Hierarchies}

\paragraph{Photon temperature multipoles.}
\begin{align}
    \Theta_{0}' &= -k\Theta_{1} - \Phi', \\
    \Theta_{1}' &= \frac{k}{3}\Theta_{0} - \frac{2k}{3}\Theta_{2}
                  + k\Psi - \dot{\tau}(\Theta_{1} - v_{b}), \\
    \Theta_{2}' &= \frac{2k}{5}\Theta_{1} - \frac{3k}{5}\Theta_{3}
                  + \dot{\tau}\left(\frac{1}{10}\Pi - \Theta_{2}\right), \\
    \Theta_{\ell}' &= \frac{k}{2\ell+1}\bigl[\ell\,\Theta_{\ell-1}
                     - (\ell+1)\Theta_{\ell+1}\bigr]
                     - \dot{\tau}\Theta_{\ell}, \qquad \ell \ge 3.
\end{align}

\paragraph{Photon polarization multipoles.}
\begin{align}
    \Theta_{0}^{(P)\prime} &= -k\Theta_{1}^{(P)} - \dot{\tau}\Theta_{0}^{(P)}, \\
    \Theta_{1}^{(P)\prime} &= \frac{k}{3}\Theta_{0}^{(P)}
                             - \frac{2k}{3}\Theta_{2}^{(P)}
                             - \dot{\tau}\Theta_{1}^{(P)}, \\
    \Theta_{2}^{(P)\prime} &= \frac{2k}{5}\Theta_{1}^{(P)}
                             - \frac{3k}{5}\Theta_{3}^{(P)}
                             + \dot{\tau}\left(\frac{1}{10}\Pi
                             - \Theta_{2}^{(P)}\right), \\
    \Theta_{\ell}^{(P)\prime} &=
    \frac{k}{2\ell+1}\bigl[\ell\,\Theta_{\ell-1}^{(P)}
    - (\ell+1)\Theta_{\ell+1}^{(P)}\bigr]
    - \dot{\tau}\Theta_{\ell}^{(P)}, \qquad \ell \ge 3.
\end{align}

\paragraph{Neutrino multipoles.}
\begin{align}
    F_{0}' &= -kF_{1} - \Phi', \\
    F_{1}' &= \frac{k}{3}F_{0} - \frac{2k}{3}F_{2} + k\Psi, \\
    F_{2}' &= \frac{2k}{5}F_{1} - \frac{3k}{5}F_{3}
             + \frac{4}{15}(\Phi' + \Psi'), \\
    F_{\ell}' &= \frac{k}{2\ell+1}\bigl[\ell\,F_{\ell-1}
               - (\ell+1)F_{\ell+1}\bigr], \qquad \ell \ge 3.
\end{align}

\paragraph{Mirror photons and neutrinos.}
Identical hierarchies with


\[
    T_{-} = xT_{+}, \qquad
    \dot{\tau} \rightarrow \dot{\tau}_{-}.
\]



\subsection{Tensor Evolution}

\paragraph{Tensor equation.}
\begin{equation}
    h'' + 2\mathcal{H}h' + k^{2}h
    = 16\pi G_{\mathrm{eff}}\,a^{2}\,\pi^{(T)}.
\end{equation}

\paragraph{Tensor anisotropic stress.}
\begin{equation}
    \pi^{(T)} = \pi^{(T)}_{\nu,+} + \pi^{(T)}_{\nu,-}.
\end{equation}

\paragraph{RDCC modulation.}
\begin{equation}
    h(k,\tau)
    = h_{0}(k,\tau)\left[1 + \alpha_{\mathrm{IR}}
      \cos\!\left(\omega_{\mathrm{mod}}\ln\frac{k}{k_{\ast}}\right)\right],
\end{equation}
with
\begin{equation}
    \omega_{\mathrm{mod}} = c_{w}\,\alpha_{\mathrm{IR}}.
\end{equation}

\subsection{Implementation Summary}

To implement RDCC in a Boltzmann code, the following modifications are required:

\begin{enumerate}
    \item Replace $G \rightarrow G_{\mathrm{eff}}$ in the Einstein equations.
    \item Add mirror-sector energy densities and pressures to the background.
    \item Add mirror-sector perturbations to the Einstein sources.
    \item Add relaxation terms $\propto \alpha_{\mathrm{IR}}$ to the continuity and Euler equations.
    \item Add mirror photon and mirror neutrino hierarchies with temperature ratio $x$.
    \item Modify the primordial tensor spectrum by the RDCC modulation factor.
\end{enumerate}

These changes are minimal, local, and controlled entirely by the single parameter
$\alpha_{\mathrm{IR}}$.

\section{Conclusion}

This companion paper has completed the final structural element required to make
Relaxation-Driven Cyclic Cosmology (RDCC) a fully testable alternative to inflationary
$\Lambda$CDM: the construction of the complete linear perturbation theory and Boltzmann
hierarchy for scalar, vector, and tensor modes in both the visible and mirror sectors.

Building on the dynamical foundation established in Companion~VI, we showed that the
relaxation parameter $\alpha(a)$ freezes extremely early and remains constant throughout all
observable epochs. As a result, all perturbative modifications of RDCC are controlled by the
single infrared parameter $\alpha_{\mathrm{IR}}$. This freezing behavior is the key to the predictive
power of RDCC: it ensures that the scalar tilt, dark-radiation excess, growth-rate deviation,
and tensor modulation are all correlated consequences of the same microscopic mechanism.

The main results of this companion paper are:

\begin{enumerate}
    \item The scalar Einstein equations acquire a modified Poisson equation with an effective
    gravitational coupling $G_{\mathrm{eff}} = G(1 + c_{g}\,\alpha_{\mathrm{IR}})$.

    \item The visible and mirror sectors contribute jointly to the gravitational potentials, while
    remaining weakly coupled through small relaxation terms proportional to
    $\alpha_{\mathrm{IR}}$.

    \item The photon and neutrino Boltzmann hierarchies retain their standard form, with RDCC
    modifications entering only through the gravitational potentials and the presence of
    mirror radiation.

    \item Mirror-sector perturbations obey parallel hierarchies with temperature ratio
    $x = T_{-}/T_{+}$ and contribute to the total anisotropic stress.

    \item Tensor perturbations inherit the distinctive RDCC signature: a primordial
    log-periodic modulation with frequency $\omega_{\mathrm{mod}} \propto \alpha_{\mathrm{IR}}$, providing
    a clean observational test in the stochastic gravitational-wave background and CMB
    $B$-modes.

    \item The full set of equations required for numerical implementation in CLASS or CAMB
    has been assembled, enabling a complete RDCC likelihood analysis.
\end{enumerate}

Together with Companions~I–VI, the present work completes the theoretical architecture of
RDCC. The framework is now:

\begin{itemize}
    \item \textbf{dynamically complete} — the evolution of $\alpha(a)$ is fully understood,
    \item \textbf{structurally complete} — the perturbation theory is fully specified,
    \item \textbf{observationally complete} — all signatures are controlled by a single infrared
    parameter,
    \item \textbf{numerically complete} — RDCC can now be implemented in Boltzmann codes.
\end{itemize}

The next step is clear. With the full perturbation system in hand, a dedicated numerical
implementation of RDCC in a Boltzmann code will allow for a complete comparison with
CMB, LSS, and GW data. This will be the subject of a future work, Companion~VIII, which
will present the first full RDCC likelihood analysis and quantify the observational viability of
the framework.

RDCC is now a fully predictive, one-parameter cosmology whose microscopic origin, global
structure, dynamical evolution, and observational consequences form a coherent and unified
whole.


\end{document}
